%
% Vorlage für Versuchsprotokolle im Grundpraktikum an der RWTH Aachen
% -------------------------------------------------------------------
%
% Bitte verwenden Sie diese Vorlage zur Erstellung Ihrer
% Protokolle.
%
% Wenn möglich, drucken Sie Ihre Protokolle bitte doppelseitig
% aus. Wenn Sie das Protokoll doppelseitig ausdrucken, verwenden Sie
% die twoside Option ...
\documentclass[twoside]{protokoll}
% ... wenn Sie nicht doppelseitig ausdrucken können, lassen Sie sie
% weg:
%\documentclass{protokoll}
\usepackage{float}
% Setzen Sie hier den jeweiligen Versuchsnamen ein:
\versuchsname{Schallgeschwindigkeit in Festkörpern}

% Setzen Sie hier Ihre(n) Namen und Ihre Matrikelnummer(n) sowie Ihre
% Gruppennummer ein:
\teilnehmer{Sophie, 466344}
\teilnehmer{Alexa, 469147}
\gruppe{C4}

\begin{document}

% Hier wird die zum Versuch gehörende Aufgaben ausgewählt, indem das
% Prozentzeichen (%) vor dem entsprechenden \input-Befehl entfernt
% wird. Die tex-Datei kann im Moodle beim jeweiligen Versuch
% heruntergeladen werden. Achten Sie darauf, dass sich alle benötigten
% Dateien im selben Verzeichnis befinden, bevor Sie diese tex-Datei
% kompilieren.
%
% Die Bearbeitung der Aufgaben findet unmittelbar in der jeweiligen
% Datei ("aufgaben_xyz.tex") statt. Dort befinden sich Umgebungen
% der Form
%
% \begin{aufgabe}
%   ... (Aufgabentext) ...
% \end{aufgabe}
%
% Hinter dem versuchsspezifischen Aufgabentext und dem "\end{aufgabe}"
% beginnen Sie jeweils Ihre Lösung der Aufgabe.
%

% Mechanik
% \input{aufgaben_pendel.tex}

% Akustik
% \section{Stehende Schallwellen in Luft}

\begin{aufgabe}{Resonanzfrequenzen mit Cassy}
  Geben Sie die Rohrlänge $L$ und ihre Unsicherheit an. Stellen Sie
  Ihre Daten als $U$-$f$-Diagramm dar: Zeichnen Sie Ihre Messpunkte
  ein und verbinden Sie diese ggfs.~mit Linien. Lesen Sie die
  Resonanzfrequenzen $f_n$ ab. Bestimmen Sie die Messunsicherheit auf
  die $f_n$ (Wiederholungsmessung, ersatzweise Abschätzung der
  Genauigkeit, mit der die Peakpositionen abgelesen werden
  können). Zeigen Sie für jede Resonanzfrequenz einen vergrößerten
  Ausschnitt Ihres Diagramms, in den Sie die gefundene Peakposition
  und ihre Unsicherheit einzeichnen. Tragen Sie die $f_n$ gegen $n$
  auf und führen Sie eine lineare Regression durch. Bestimmen Sie die
  Schallgeschwindigkeit und ihre Unsicherheit.
\end{aufgabe}

First we measured the length of the tube we used in our experiment, we obtained the following results conducting independent experiments:
\vspace{1em}
\begin{table}[H]
  \begin{tabular}{c|cc|l}
    Length & Measurements & & Mean \\ \hline
    L & $41,95$ cm & $42$ cm & $41,975$ cm 
  \end{tabular}
\end{table}
we then calculate the uncertainty in our length measurement by quadratically adding our stadard deviation and the sys. error given by the manufacturer of the tool of measurement (in this case a measuring tape with 0,7 mm uncertainty):
$$
\sigma _{L}  =  \sqrt{ \sigma^{2} + \sigma _{sys}^{2} }
$$
$$
\sigma _{L}  =  \sqrt{ 0,35^{2} + (\frac{0,07}{\sqrt{ 3 }})^{2} } \text{ mm} = 0,357 \text{ mm}
$$



we then used scripts to find the first five resonance frequencies, we then sorted our data and noticed that in some of our Measurements we had more than one voltage corresponding to a frequency value so we averaged them and took the peaks of the resulting plot for clarity.
for each of our peaks, we decided that an uncertainty of $2$ Hz is enough to contain the top $10\%-20\%$ of our peak length, and so we took $2$ Hz as conservative frequency uncertainty (highlighted with a red error bar). We then plot our found frequencies and their uncertainties in dedicated plots:




\begin{figure}[H]  
    \centering
    \includegraphics[width=0.8\textwidth]{f_u.pdf}  
    \caption{plot for frequency and voltage}  
    \label{fig:f_1}  
\end{figure}
\begin{figure}[H]  
    \centering
    \includegraphics[width=0.8\textwidth]{f_1.pdf}  
    \caption{first frequency}  
    \label{fig:f_1}  
\end{figure}
\begin{figure}[H]  
    \centering
    \includegraphics[width=0.8\textwidth]{f_2.pdf}  
    \caption{second frequency}  
    \label{fig:f_2}  
\end{figure}
\begin{figure}[H]  
    \centering
    \includegraphics[width=0.8\textwidth]{f_3.pdf}  
    \caption{third frequency}  
    \label{fig:f_3}  
\end{figure}
\begin{figure}[H]  
    \centering
    \includegraphics[width=0.8\textwidth]{f_4.pdf}  
    \caption{fourth frequency}  
    \label{fig:f_1}  
\end{figure}
\begin{figure}[H]  
    \centering
    \includegraphics[width=0.8\textwidth]{f_5.pdf}  
    \caption{fifth frequency}  
    \label{fig:f_1}  
\end{figure}

we then give the results of our maximum peaks as follows:
\vspace{1em}
\begin{itemize}
  \item $413$ Hz: first resonance frequency.
  \item $818$ Hz: second resonance frequency.
  \item $1225$ Hz: third resonance frequency.
  \item $1632$ Hz: fourth resonance frequency.
  \item $2040$ Hz: fifth resonance frequency.
\end{itemize}

\vspace{1em}

Then through a linear regression we fit our resonance frequencies to their number, we obtain the following result:

\begin{verbatim}
  R = (406.8+/-0.6) Hz,
  b = (5.2+/-2.1) Hz
  chi2/dof = 0.7 / 3  
  corr = -0.904535
\end{verbatim}
\begin{figure}[H]  
    \centering
    \includegraphics[width=0.8\textwidth]{reg_f_1.pdf}  
    \caption{linear regression fit for frequencies and their numbers}  
    \label{fig:f_1}  
\end{figure}

We then combined our frequencies divided by their number in a variable $f_n /n$ and calculated its uncertainty through weighted average of the resulting division, overall we obtained the following variables and their uncertainties, which leads us to the speed of sound:
$$
v = \frac{f_{n}}{n} \cdot 2 L
$$
\begin{verbatim}
  L: 41,975 +/- 0,357 cm
  f_n / n : 408,2 +/- 0,27 Hz
  v : 342,7 +/- 0,37 m/s
\end{verbatim}














% Bitte belassen Sie die Aufgabentexte in Ihrem Protokoll und beginnen
% Sie hier mit der Lösung der ersten Aufgabe:

% Achten Sie darauf, alle verwendeten Formeln anzugeben und die
% verwendeten Formelzeichen zu definieren.










\begin{aufgabe}{Resonanzfrequenzen mit dem Oszilloskop}
  Tabellieren Sie die gefundenen Resonanzfrequenzen $f_n$. Zur
  Bestimmung der Messunsicherheit auf die $f_n$ führen Sie
  Wiederholungsmessungen durch oder schätzen Sie die Genauigkeit ab,
  mit der Sie die Resonanzfrequenzen mit dem Oszilloskop finden
  können. Bestimmen Sie aus Ihren Messdaten die Schallgeschwindigkeit
  und ihre Unsicherheit. Tragen Sie dazu die $f_n$ gegen $n$ auf und
  führen Sie eine lineare Regression durch.
\end{aufgabe}







\begin{aufgabe}{Vermessung der stehenden Welle}
  Tabellieren Sie die gefundenen Positionen $x_n$ der Druckknoten und
  -bäuche. Geben Sie die jeweils verwendete Resonanzfrequenz $f$ und
  ihre Unsicherheit an. Zur Bestimmung der Messunsicherheit auf die
  $x_n$ führen Sie Wiederholungsmessungen durch oder schätzen Sie die
  Genauigkeit ab, mit der Sie die Druckknoten bzw. -bäuche durch
  Verschiebung des Mikrofons und Beobachtung der Amplitude am
  Oszilloskop einstellen können. Tragen Sie die Positionen $x_n$ gegen
  $n$ auf und führen Sie eine geeignete lineare Regression
  durch. Bestimmen Sie aus Ihren Ergebnissen die Schallgeschwindigkeit
  und ihre Unsicherheit. Berechnen Sie den gewichteten Mittelwert über
  die Ergebnisse aus Ihren Messreihen bei verschiedenen
  Resonanzfrequenzen.
\end{aufgabe}

\begin{aufgabe}{Zusammenfassung und Vergleich}
  Stellen Sie die Ergebnisse aller drei Teilversuche
  nebeneinander. Berechnen Sie ggfs. ein gewichtetes Mittel der
  Schallgeschwindigkeit unter Berücksichtigung der Korrelationen der
  Unsicherheiten. Vergleichen Sie Ihr Ergebnis mit dem Literaturwert
  für die gemessene Raumtemperatur.
\end{aufgabe}

\section{Schallgeschwindigkeit in Festkörpern}

\begin{aufgabe}{Rohdaten}
  Stellen Sie die gemessenen Werte von $L$, $D$ und $M$ in
  tabellarischer Form dar. Visualisieren Sie für jeden vermessenen
  Metallstab einen typischen Schwingungsverlauf sowie dessen
  Fourierspektrum in geeigneter Weise.
\end{aufgabe}

% Bitte belassen Sie die Aufgabentexte in Ihrem Protokoll und beginnen
% Sie hier mit der Lösung der ersten Aufgabe:

% Achten Sie darauf, alle verwendeten Formeln anzugeben und die
% verwendeten Formelzeichen zu definieren.

\begin{aufgabe}{Auswertung}
  Bestimmen Sie aus den aufgezeichneten Schwingungsvorgängen die
  Schwingungsfrequenzen und tabellieren Sie die erhaltenen
  Werte. Ermitteln Sie die Streuung der Frequenzwerte und bestimmen
  Sie daraus den statistischen Fehler auf den Mittelwert. Erläutern
  und illustrieren Sie Ihr Vorgehen zur Bestimmung der systematischen
  Unsicherheiten auf die Frequenz. Diskutieren Sie die Unsicherheiten
  auf $L$, $D$ und $M$. Berechnen Sie den E-Modul $E$ und die Dichte
  $\rho$ der vorliegenden Stangen und die zugehörigen statistischen
  und systematischen Messunsicherheiten. Fassen Sie Ihre Ergebnisse
  tabellarisch zusammen. Diskutieren Sie, welche Fehlerbeiträge den
  Gesamtfehler auf $E$ dominieren. Vergleichen Sie Ihre Ergebnisse für
  $E$ und $\rho$ mit einschlägigen Literaturwerten.
\end{aufgabe}





% Thermodynamik
% \section{Adiabatenkoeffizient der Luft}

\begin{aufgabe}{Analyse der Kugelschwingungen}
  Beschreiben Sie kurz Ihr Vorgehen bei der Messung. Zeigen Sie
  exemplarisch für jeden Erlenmeyerkolben sowie die große Mariottesche
  Flasche je einen Schwingungsverlauf und seine Auswertung. Tabellieren
  Sie Ihre aus allen Schwingungsverläufen ermittelten Ergebnisse und
  die zugehörigen Kugelpositionen im Rohr.
\end{aufgabe}

% Bitte belassen Sie die Aufgabentexte in Ihrem Protokoll und beginnen
% Sie hier mit der Lösung der ersten Aufgabe:

% Achten Sie darauf, alle verwendeten Formeln anzugeben und die
% verwendeten Formelzeichen zu definieren.
\section{Overview and rudemintary measurements}

We first fix the flask in place and attach the tubes tightly to make sure no air escapes. we then connect the glass tube to the top of the flask and make sure it is vertical. we make sure the pressure sensor is connected and the whole setup is ready to conduct the experiment, we then lower the ball to a position where fall damage is not a concern using the magnets. once we reach a suitable spot, we start conducting the experiment by releasing the ball and recording the pressure change until the ball reaches the equilibrium position, we then measure the hight of the ball using the ruler and repeat the experiment for few times.

\vspace{1em}
we summarise the measurements in the following tables:
\begin{table}[h]
  \begin{tabular}{l|ccc}
    Properties & \multicolumn{3}{|c}{Measurements } \\ \\ \hline
     m (g) & 16,5 & 16,4 & 16,4  \\
           &  &  &  \\ \hline
     p (hPa)& 981 & 980 &  \\
           &  &  &  \\ \hline
     D (mm) & 15,98 &15,55  &  \\
           &  &  &  \\ \hline
     T$_{0}$ (s) &0,38  & 0,39 & 0,38 \\
           & 0,39 &  0,38 &0,39  \\ \hline
  \end{tabular}
  \caption{Measurements for the 1 Liter flask.}
  \label{tab:beispieltabelle}
\end{table}

\begin{itemize}
  \item We measured the mass of the ball using the scale available in the lab.
  \item We measured the pressure using the pressure sensor in lab along the first few minutes.
  \item We measured the diameter of the ball using the micrometer caliper available in the lab.
  \item We measured the period of the oscillation using the Cassy measurements and the time difference between multiple peaks.
\end{itemize}

\section{Errors and assumptions}

\begin{table}[h]
  \centering
  \begin{tabular}{l|ccc}
    Measurement tool & Stat. error  & Sys. error & Combined errors\\ \\ \hline
     Scale (g) &  &  &  \\
           &  &  &  \\ \hline
     Pressure sensor (hPa)&  &  &  \\
           &  &  &  \\ \hline
     Micrometer caliper (mm) &  &  &  \\
           &  &  &  \\ \hline
     Cassy (s) &  &  &  \\
           &  &  &  \\ \hline
  \end{tabular}
  \caption{Beispieltabelle.}
  \label{tab:beispieltabelle}
\end{table}







\begin{aufgabe}{Messungen der übrigen Observablen}
  Beschreiben Sie die Messungen der übrigen Observablen. Geben Sie
  jeweils Ihre Mess\-er\-geb\-nis\-se an und bestimmen Sie die
  zugehörigen Messunsicherheiten. Verwenden Sie ggfs. Tabellen, um
  Ihre Ergebnisse übersichtlich darzustellen.
\end{aufgabe}









\begin{aufgabe}{Ergebnis und Fehlerrechnung}
  Berechnen Sie $\kappa$ für die einzelnen Messungen. Beschreiben
  Sie, wie Sie die Unsicherheit auf den Adiabatenkoeffizienten
  $\kappa$ bestimmen. Tabellieren Sie exemplarisch für die große
  Mariottesche Flasche die Größe der einzelnen Beiträge zur
  Unsicherheit auf $\kappa$. Geben Sie dabei jeweils an, ob eine
  Unsicherheit von Messung zu Messung unkorreliert ist, oder ob es
  sich um eine korrelierte Unsicherheit handelt. Geben Sie bei den
  korrelierten Unsicherheiten an, ob sie nur für ein einzelnes
  Gasvolumen oder über alle Messungen hinweg korreliert sind.
  
  Stellen Sie die für die einzelnen Messungen erhaltenen Werte von
  $\kappa$ in einem Plot dar. Machen Sie die von Messung zu Messung
  unkorrelierten Fehler jeweils durch einen Fehlerbalken kenntlich und
  die (geeignet zusammengefassten) korrelierten Fehler durch
  Bänder. Geben Sie Ihr Endergebnis für $\kappa$ an. Vergleichen Sie
  mit dem Literaturwert und diskutieren Sie Ihr Ergebnis.
\end{aufgabe}


\begin{table}[h]
  \centering
  \begin{tabular}{l|ccc}
    Properties & \multicolumn{3}{|c}{Measurements } \\ \\ \hline
     M (g) &  &  &  \\
           &  &  &  \\ \hline
     p(hPa)&  &  &  \\
           &  &  &  \\ \hline
     D(mm) &  &  &  \\
           &  &  &  \\ \hline
  \end{tabular}
  \caption{Beispieltabelle.}
  \label{tab:beispieltabelle}
\end{table}
% \input{aufgaben_dampfdruckkurve.tex}

% Optik
% \input{aufgaben_prismenspektrometer.tex}
% \section{Michelson-Interferometer}

\begin{aufgabe}{Rauschmessung}
  Tabellieren Sie die Rohdaten Ihrer Messungen zur wiederholten
  Einstellung und Ablesung an einem festen Maximum. Berechnen Sie aus
  ihren Messwerten die empirische Standardabweichung als Maß für die
  Messunsicherheit bei den folgenden Messungen mit dem
  Feinsteinstelltrieb.
\end{aufgabe}

% Bitte belassen Sie die Aufgabentexte in Ihrem Protokoll und beginnen
% Sie hier mit der Lösung der ersten Aufgabe:

% Achten Sie darauf, alle verwendeten Formeln anzugeben und die
% verwendeten Formelzeichen zu definieren.

\subsection{Noise level through repeated measurements.}
\textbf{Goal:} \\
The detrermination of the noise through the measurement of the standard deviation of repeated length measurements with a fixed starting point. We also need to mention that we were allowed by the instructor to estimate the fractional values of the reading on the caliper, which is why we have values such as $7.419$ mm in our measurements, which is an estimation of the reading on the caliper.
\vspace{1em}


\textbf{Methodology:}
\begin{itemize}
    \item We perform 8 measurements of the length corresponding to an difference of number of maxima of $\Delta m=10$.
    \item We start by noting our starting position at $s_0 = 7.485$ mm where we locate a maximum. we then move the micrometer screw until we observe $\Delta m = 10$ maxima and note the new position $s_i$ for each measurement.
    \item We then calculate the length difference $\Delta s_i = s_i - s_0$ for each measurement and use these values to calculate the arithmetic mean and standard deviation. we then use it as a measure of the noise in our measurements.
\end{itemize}

\textbf{Data:}

We then tabulate our measurements as follows, making sure to include the calculated length differences $\Delta s_i$ for each measurement:

\begin{table}[H]
    \centering
    \label{tab:messwerte}
    \begin{tabular}{ccc}
        \toprule
        reading & value (mm) & $\Delta s_i =  s_i - s_0$ (mm)\\
        \midrule
        $s_1$ & 7.419 & 0.066\\
        $s_2$ & 7.420 & 0.065\\
        $s_3$ & 7.412 & 0.073\\
        $s_4$ & 7.418 & 0.067\\
        $s_5$ & 7.415 & 0.070\\
        $s_6$ & 7.410 & 0.075\\
        $s_7$ & 7.412 & 0.073\\
        $s_8$ & 7.411 & 0.074\\
        \bottomrule
    \end{tabular}
    \caption{Measurements for the $m=10$ readings on the caliper}
\end{table}

\textbf{Evaluation:} \\
After obtaining our data, we then calculated the arithmetic mean and the empirical standard deviation, and the error on the mean as follows:
\begin{align}
    \bar{s} &= \frac{1}{8} \sum_{k=1}^{8} \Delta s_{k}, \\
    s_{\Delta s} &= \sqrt{\frac{1}{8-1} \sum_{k=1}^{8} \left(\Delta s_k - \bar{s}\right)^2}, \\
    \Delta \bar{s} &= \frac{s_{\Delta s}}{\sqrt{8}}.
    \label{eq:uncertainty_workflow for measurements}
\end{align}

\begin{verbatim}
    arithmetic mean: 0.070 mm
    empirical standard deviation: 0.0039 mm
    error on the mean: 0.0013 mm
    end result: 0.070 +/- 0.0013 mm
\end{verbatim}

\textbf{comments:} \\
through our last measurement, we obtained a value of $ \bar{s}  = 0.070 \pm 0.0013$ mm for the length difference corresponding to $\Delta m = 10$ maxima. This means that the noise in our measurements is equal to  $0.0013$ millimeters, which is a reasonable value for this type of measurement.



























\begin{aufgabe}{Kalibration des Übersetzungsfaktors}
  Bestimmen Sie den Übersetzungsfaktor $k$ des Feinsteinstelltriebes.
  Verwenden Sie dazu mehrere Messreihen über einen kürzeren Bereich
  der Skala $s$ der Mikrometerschraube oder eine Messreihe über einen
  längeren Bereich. Tabellieren Sie $s$ in Abhängigkeit der Anzahl der
  durchlaufenen Maxima $m$ und führen Sie geeignete lineare
  Regressionen zur Bestimmung von $k$ durch. Unter Umständen müssen
  Sie den Bereich der Skala dabei unterteilen und jeweils verschiedene
  Werte von $k$ bestimmen. Geben Sie Ihren Wert bzw.~Ihre Werte für
  $k$ samt Messunsicherheit an.
\end{aufgabe}


\begin{aufgabe}{Vermessung der Wellenlänge des Nd:YV04-Lasers}
  Bestimmen Sie die Wellenlänge $\lambda$ des
  Nd:YV04-Lasers. Tabellieren Sie $s$ in Abhängigkeit der Anzahl der
  durchlaufenen Maxima $m$ und führen Sie lineare Regressionen
  durch. Tabellieren Sie die Ergebnisse Ihrer
  Anpassungsrechnungen. Bestimmen Sie in geeigneter Weise ein
  Endergebnis für $\lambda$ und die zugehörige
  Messunsicherheit. Diskutieren Sie die Natur der Beiträge zur
  Gesamtunsicherheit. Vergleichen Sie Ihr Endergebnis mit der
  Herstellerangabe.
\end{aufgabe}





% Elektrizitätslehre
% \input{aufgaben_lcr_serienschwingkreis.tex}

\end{document}
